\documentclass[a4paper,12pt]{article}
\usepackage[british]{babel}
\usepackage[style=apa,backend=biber]{biblatex}
\usepackage{csquotes}
\usepackage{fancyhdr}
\usepackage[margin=1in]{geometry}
\usepackage[hidelinks,linkcolor=blue,citecolor=blue,urlcolor=blue,colorlinks]{hyperref}
\usepackage{indentfirst}
\usepackage{newtxtext}
\usepackage{newtxmath}
\usepackage[document]{ragged2e}
\usepackage{setspace}
\usepackage[explicit]{titlesec}

\addbibresource{references.bib}

\doublespacing

% Paragraph formatting (APA7: 0.5in indent, no extra paragraph spacing)
\setlength{\RaggedRightParindent}{0.5in}
\setlength{\parindent}{0.5in}
\setlength{\parskip}{0pt}

% Running head (all pages including title page)
\pagestyle{fancy}
\fancyhf{}
\fancyhead[L]{\runninghead}
\fancyhead[R]{\thepage}
\renewcommand{\headrulewidth}{0pt}

% APA7 heading levels
% Level 1: Centred, Bold
\titleformat{\section}
    {\normalfont\bfseries\centering}{}{0pt}{#1}
\titlespacing*{\section}{0pt}{0.3\baselineskip}{0pt}

% Level 2: Flush Left, Bold
\titleformat{\subsection}
    {\normalfont\bfseries}{}{0pt}{#1}
\titlespacing*{\subsection}{0pt}{0.3\baselineskip}{0pt}

% Level 3: Flush Left, Bold Italic
\titleformat{\subsubsection}
    {\normalfont\bfseries\itshape}{}{0pt}{#1}
\titlespacing*{\subsubsection}{0pt}{0.3\baselineskip}{0pt}

% Level 4: Indented, Bold, Ending With a Period. Text continues on same line.
\titleformat{\paragraph}[runin]
    {\normalfont\bfseries}{}{0pt}{#1.}
\titlespacing*{\paragraph}{\parindent}{0.3\baselineskip}{\wordsep}

% Level 5: Indented, Bold Italic, Ending With a Period. Text continues on same line.
\titleformat{\subparagraph}[runin]
    {\normalfont\bfseries\itshape}{}{0pt}{#1.}
\titlespacing*{\subparagraph}{\parindent}{0.3\baselineskip}{\wordsep}

% Title page fields
\newcommand{\runninghead}{}
\newcommand{\essaytitle}{}
\newcommand{\essayauthor}{}
\newcommand{\essaycourse}{}
\newcommand{\essayinstitution}{}
\newcommand{\essayinstructor}{}
\newcommand{\essaydate}{}

% APA7 title page
\newcommand{\apatitlepage}{%
    \thispagestyle{fancy}
    \begin{center}
        \vspace*{3\baselineskip}
        {\bfseries\essaytitle}\\[\baselineskip]
        \essayauthor\\
        \essaycourse\\
        \essayinstitution\\
        \essayinstructor\\
        \essaydate
    \end{center}
    \newpage
}

% De facto Level 1 heading: paper title repeated at start of body text (APA7)
\newcommand{\apatextheading}{%
    \begin{center}
        {\bfseries\essaytitle}
    \end{center}
}

% formatting shortcuts
\newcommand{\I}[1]{\textit{#1}}
\newcommand{\B}[1]{\textbf{#1}}

% Commands for secondary citations
\newcommand{\textseccite}[3]{#1 (#2, as cited in \nptextcite{#3})}

\newcommand{\possessivecite}[1]{\citeauthor{#1}'s \parencite*{#1}}

\renewcommand{\runninghead}{TOOL OR WEAPON}
\renewcommand{\essaytitle}{A Tool or a Weapon? Analyzing the Exploitation of Corporate Jargon in Workplace Communication}
\renewcommand{\essayauthor}{Zhan Ho Jacob Shing}
\renewcommand{\essaycourse}{CCHU9094: From Grunts to Grammar: The Evolution of Human Communication}
\renewcommand{\essayinstitution}{The University of Hong Kong}
\renewcommand{\essayinstructor}{Dr. Manwa Lawrence Ng}
\renewcommand{\essaydate}{13 April 2026}

\begin{document}

\apatitlepage

Consider a scenario set in a quarterly business review meeting at a middle-sized corporation that has recently fallen short of its financial targets substantially.
The participants include a senior executive, a mid-level manager who has been in the firm for several years, a junior employee who has been with the company for less than a year, and several other staff members.
The situation of the firm is beyond recovery and downsizing is inevitable.

The ultimate goal of the senior executive is not a simple disclosure of the situation and the decision of layoffs, but rather using a strategic communication approach, or as we call it, ``corporate jargon'', to prevent panic, maintain morale, preserve a facade of control and competence, and motivate and reassure the remaining workforce.
In simpler terms, to control and to heal.

This essay argues that corporate jargon is a series of meticulous linguistic choices \parencite{Zamfir2024} that function primarily as tools of social control and manipulation through the mechanisms of euphemism, obfuscation through syntactic manoeuvering, and shibboleth. The essay further argues that while the use of corporate jargon is intended to be a tool for constructive communication, it can easily become a weapon that alienates different audiences.

\section{The Linguistic Mechanics of Obfuscation}

This section explores the verbal strategies employed in the corporate jargon used by the senior executive, illustrated by examples for the scenario described above.

\subsection{A Linguistic Choice: Euphemism and Euphemism Treadmill}

\textcite{Stein1999} defines euphemism as a substituted fragment of language that conveys less negative connotations than the original term, which is often considered offensive or harsh.
\textcite{Rittenburg2016} see euphemism as a way of reducing transparency, thereby serving the function of obfuscation.
In other words, euphemisms are used as a semantic softening agent.
For the scenario of this essay, the senior executive may employ terminologies such as, as \textcite{Vickers2002} suggested, ``right-sizing'', ``downsizing'', or ``shifted resources'' instead of the more straightforward terms ``layoffs'', if not ``firings''.
As for the substantial shortfall in meeting targets, rather than a disaster-like ``failure'', one may frame it as a ``challenge'' or a ``need for calibration of expectations''.
The use of euphemisms in this context often shifts the focus from the negative aspects of the situations to a more positive connotation, which can help maintain morale while also obscuring responsibility, achieving the purposes of control (managing the emotions and response of the employees) and healing (preserving the image of the executive as a competent leader).

However, one may argue that the controlling and healing effects from the perpetual use of these terms gradually wear off, as the audience starts to map the euphemisms to the underlying negative realities.
Indeed, as \textcite{LumenLearning2026} suggested, in the modern corporate life, the term ``downsizing'' is generally considered to be delivering a negative connotation.
Nevertheless, this is an expected outcome of the use of euphemisms, as \textcite{Stein1999} also noted that euphemisms draws away attention from the negative ideas in one level, but they, in a way, also imply what is being euphemised.

To overcome the weakened euphemisms, new euphemisms need to be created, which is known as the ``euphemism treadmill'' \parencite{Pinker1994}.
On this treadmill, the words are constantly renewed, powering the mechanics of euphemism to sustain itself over decades of corporate communication.
While more senior or experienced employees may be already familiar with decoding phrases like ``downsizing'', the change to ``right-sizing'' or ``workforce optimisation'' may buy the executive some time and further defer the negative connotations for both senior and junior employees.

In \possessivecite{Jakobson1960} terms, euphemisms of corporate jargon are strategic exploitation of the expressive and phatic functions (to inform and to keep social bonds) of language, while also skilfully mutating the referential function (to refer to the underlying ugliness) to achieve the purposes of control and healing.

\subsection{Syntactic Manoeuvring: Nominalisation and Passive Voice}

While euphemism lives on the lexical level, corporate jargon also exploits the syntactic or grammatical level to achieve obfuscation.
It is not uncommon for corporate jargon to intentionally use complex syntactic structures to create further obfuscation and blur responsibility.
Whereas \textcite{Halliday1985} view nominalisation as an essential tool for English, \textcite{Bello2016} examines its cognitive implications -- turning a dynamic action, often acted by a person, into a static, reified object.
This linguistic act diverts humans' minds into isolating the objects from the agent and the person being acted upon \parencite{Barsalou2010}, and thus creating a sense of detachment of responsibility and agency.

For example, rather than saying ``the company has failed to meet its targets'', the senior executive could have instead phrased it as ``a recent challenge has been noted in the area of target achievement'', which tranforms the verb ``to fail'' and ``to meet'' into ``challenge'' and ``achievement'', respectively.
As for the layoffs, instead of saying ``we have decided to lay off 20\% of the workforce'', it could have been ``a decision of workforce simplification has been made''.

One should immediately notice that in the above examples, nominalisation is often accompanied by the employment of passive voice.
Indeed, the passive construction is a common verbal strategy for drawing attention towards the action and the result, rather than the agents \parencite{Memmedli2025}.
In such constructions, the agent is introduced by a ``by'' clause, which can further be omitted to entirely erase agency, and in turn, responsibility \parencite{Bailey1997}.

Not coincidentally, the use of nominalisation and passive voice serves the same purpose of obfuscation through detachment of agency and responsibility.
This sort of language manipulation is particularly effective in the setting of corporate communication through delegating the responsibility to virtually no one, while the senior executives remain in control of the narrative.
The objectification of the actions taken also eliminates the tied sentiments of guilt and blame, which is crucial for maintaining morale, and hence, the healing function of language.

\subsection{The Shibboleth Effect}

Having explored the syntactic linguistic strategies, this essay now dives deeper into the sociolinguistic mechanism of shibboleth.
Served as a insider/outsider test historically \parencite{McNamara2005}, the shibboleth effect has now been established as a linguistic marker and tool of social differentiation and bifurcation \parencite{Busch2021}.

\section{Outcomes}

\subsection{Intended Outcomes}

\subsection{Unintended Outcomes}

\section{Reflections and Conclusion}

\clearpage

\printbibliography

\end{document}
